\documentclass{article}
\usepackage{amsmath}
\usepackage{amsfonts}
\usepackage{graphicx}
\usepackage{wrapfig}
\usepackage{hyperref}
\usepackage{minted}
\usepackage{xcolor}
\usepackage{caption}

\graphicspath{{build/media/}}
\definecolor{grey}{rgb}{.95, .95, .95}
\definecolor{mygreen}{rgb}{.6 .95 .1}
\pagecolor{yellow!4}

\title{Introduction to computing for physicists}
\author{Your dearest, green}

\begin{document}
\maketitle

\pagebreak

\section{Introduction}
\paragraph{Installation}
Create a file on your computer named ``main.c'',
inside your are going to write the following:
\begin{minted}[bgcolor=mygreen!10]{c}
  #include <stdio.h>

  int main(void) {
    printf("Hello Seamen!\n");
    return 0;
  }
\end{minted}
now open a \textit{terminal/shell} of your choice
to \textit{compile} and \textit{run} the ``main.c''
file with the following:
\begin{minted}[bgcolor=mygreen!10]{bash}
  # I am on Linux and use 'gcc' depending
  # on your operating system you may need
  # to use something else.
  gcc -o main.c main
  ./main
\end{minted}
if the phrase ``Hello Seamen!'' appeared on your
terminal its good, you can move to the next paragraph.
If nothing happened make sure everything has been properly
written, if an error occurs \textit{google} ``How to install
C compiler on Windows/Mac/Linux'' or ask a trusted friend to
do so for you. Repeat this process until the
desired result has been reached.

\paragraph{Pre-requisites}
In order to move to the next chapter it is required
to understand the behaviour of the following program
without any external help:
\begin{minted}[bgcolor=mygreen!10]{c}
  // 'int' is an abbreviation for integer
  // '=' is an assignment, translate it with 'set to'
  int factorial(int n) {
    int x = 1;
    loop:
    if (n > 0) {
      x = x * n;
      n = n - 1;
      goto loop;
    }
    else {
      return x;
    }
  }
\end{minted}
Can you guess the \textit{return} value of $x$ if $n = 5$?
If you understood the code correctly it should be $x = 120$.


\pagebreak

\section{Program representation}

\subsection{Flow diagram representation}
The rules for the representation of a flow diagram are the following,
right after, a familiar diagram is shown.
\begin{itemize}
\item Circles indicate insertion and termination points,
  respectively named ``start'' and ``end''.
\item Slanted parallelograms indicate insertion variables
  also known as \textbf{input} and termination variables also
  know as \textbf{output}.
\item Rhombi are \textbf{conditionals} their \underline{exit flow is not fixed},
  its a function of their inner state.
  For our purposes the possible exit flows are \textit{two}, one if the condition
  \underline{is} satisfied and another one if it \underline{is not}.
\item Rectangles are \textbf{procedures}, their \underline{exit flow is fixed}.
\end{itemize}

\begin{figure}[h]
  \centering
  \includegraphics[width=\textwidth]{flow-diagram-example.png}
  \caption{
    This is the flow diagram of the ``factorial'' program
    previously shown.
  }
  \label{fig:flow-diagram-example}
\end{figure}

\pagebreak

\subsection{Procedural representation}
The ordering given by arrows is now replaced 
There are no more arrows, each line is executed in order
from top to bottom up until an \textit{indirection} command
such as ``GOTO'' which diverts the next execution to a
special line marked with a ``FLAG'' statement. This is a
typical scheme for how computers execute code.
\begin{minted}[bgcolor=mygreen!10]{text}
  START n
  SETTO x 1
  FLAG LOOP
  IFSTART GREATERTHAN n 0
  SETTO x PRODUCT x n
  SETTO n SUBTRACTION n 1
  GOTO LOOP
  IFSTOP
  ELSESTART
  STOP x
  ELSESTOP
\end{minted}

\pagebreak

\subsection{Flow patterns}
The flow of a program is a very important aspect to understand its behaviour,
all useful programs have a source-sink aspect to them, all arrows begin from
a source and they must all end in a sink at some point.

\noindent
\begin{minipage}{0.5\textwidth}
  \vspace{1cm}
  \includegraphics[width=\textwidth]{flow-if-else.png}
  \captionof{figure}{This is an \textbf{if-else} pattern, it starts from $C$.}
  \label{fig:flow-if-else}
\end{minipage}
\hspace{0.05\textwidth}
\begin{minipage}{0.4\textwidth}
  The \textit{if-else} pattern \underline{conserves} the tree structure of the program.
  If $1$ and $2$ are halting programs than the whole tree will halt.
  \begin{minted}[bgcolor=mygreen!10]{c}
    if (C) {
      // 1
    }
    else {
      // 2
    }
  \end{minted}
\end{minipage}

\noindent
\begin{minipage}{0.4\textwidth}
  \vspace{1cm}
  The \textit{if-elseif} pattern is a pattern extension of the \textit{if-else} pattern.
  It \underline{conserves} the tree structure of the program.
  If $1 \cdots n$ are halting programs than the whole tree will halt.
  \begin{minted}[bgcolor=mygreen!10]{c}
    if (1) {
      // 1
    }
    else if (...) {
      // ...
    }
    else {
      // N
    }
  \end{minted}
\end{minipage}
\hspace{0.05\textwidth}
\begin{minipage}{0.5\textwidth}
  \includegraphics[width=\textwidth]{flow-if-elseif.png}
  \captionof{figure}{This is an \textbf{if-elseif} pattern, it starts from $C$.}
  \label{fig:flow-if-elseif}
\end{minipage}

\noindent
\begin{minipage}{0.5\textwidth}
  \vspace{1cm}
  \includegraphics[width=\textwidth]{flow-if.png}
  \captionof{figure}{This is a free \textbf{if} pattern, it starts from $C$.}
  \label{fig:flow-if}
\end{minipage}
\hspace{0.05\textwidth}
\begin{minipage}{0.4\textwidth}
  The \textit{free if} pattern is a \underline{reduction} of an \textit{if-else} pattern.
  If $1 \otimes 2$ and $2$ halts than the whole tree will halt
  \begin{minted}[bgcolor=mygreen!10]{c}
    if (C) {
      // 1
    }
    // 2
  \end{minted}
\end{minipage}

\noindent
\begin{minipage}{0.4\textwidth}
  \vspace{1cm}
  The \textit{while} pattern is guaranteed to halt
  if the condition $C$ has an \textit{upper-bound}
  check as a logical conjunction and is a function
  of a \textit{strictly increasing} value (the same
  is true for a lower-bound with a stricly decreasing
  value).
  \begin{minted}[bgcolor=mygreen!10]{c}
    while (C) {
      // B
    }
  \end{minted}
\end{minipage}
\hspace{0.05\textwidth}
\begin{minipage}{0.5\textwidth}
  \includegraphics[width=\textwidth]{flow-while.png}
  \captionof{figure}{This is \textbf{while} pattern, it starts from $C$.}
  \label{fig:flow-while}
\end{minipage}

\noindent
\begin{minipage}{0.5\textwidth}
  \vspace{1cm}
  \includegraphics[width=\textwidth]{flow-do-while.png}
  \captionof{figure}{This is a \textbf{do-while} pattern, it starts from $C$.}
  \label{fig:flow-do-while}
\end{minipage}
\hspace{0.05\textwidth}
\begin{minipage}{0.4\textwidth}
  The \textit{do-while} pattern exhibits the same
  halting behaviour as the \textit{while} pattern;
  the difference is that the \textit{body} of the
  loop $B$ will always resolve \underline{at least}
  once.
  \begin{minted}[bgcolor=mygreen!10]{c}
    do {
      // B
    } while (C);
  \end{minted}
\end{minipage}

\noindent
\begin{minipage}{0.4\textwidth}
  \vspace{1cm}
  The \textit{for} pattern exhibits the same
  halting behaviour as the \textit{while} pattern;
  in addition if $C \otimes B = C \oplus B$ then
  the strictly increasing/decreasing value is
  controlled by $IT$.
  \begin{minted}[bgcolor=mygreen!10]{c}
    for (IN; C; IT) {
      // B
    }
  \end{minted}
\end{minipage}
\hspace{0.05\textwidth}
\begin{minipage}{0.5\textwidth}
  \includegraphics[width=\textwidth]{flow-for.png}
  \captionof{figure}{This is \textbf{for} pattern, it starts from $IN$.}
  \label{fig:flow-for}
\end{minipage}

\paragraph{Reductions}


\pagebreak

\subsection{Exercises}
For the following exercises read the instructions, draw the flow diagram,
apply reductions if needed and implement the solution in C syntax.

\paragraph{Scale a row}
Given a row of floats $R[n]$ of size $n$ and a factor $fac$, implement a program
to scale each column entry by the same factor (row scaling)
\begin{minted}[bgcolor=mygreen!10]{c}
  void row_scale(int n, float R[n], float fac);
\end{minted}

\paragraph{Add a row}
Given a two rows of floats $R0[n]$ and $R1[n]$ of size $n$ and a factor $fac$,
implement a program that adds together adjacent columns of $R0[n]$ and $R1[n]$
the latter being scaled by a factor of $fac$.
\begin{minted}[bgcolor=mygreen!10]{c}
  void row_add(int n, float R0[n], float R1[n], float fac);
\end{minted}

\paragraph{Swap a row}

\paragraph{Print a matrix}
Given a matrix $M[m][n]$ of size $m \times n$, implement a program that prints
the matrix onto the screen.
\begin{minted}[bgcolor=mygreen!10]{c}
  void matrix_print(int m, int n, float M[m][n]);
\end{minted}
keep in mind that to print something like
\begin{minted}[bgcolor=mygreen!10]{text}
  [1.00 6.67]
  [3.14 6.96]
\end{minted}
the code is the following
\begin{minted}[bgcolor=mygreen!10]{c}
  printf("[%.2f ", 1.00);
  printf("%.2f]\n", 6.67);
  printf("[%.2f ", 3.14);
  printf("%.2f]\n", 6.69);
\end{minted}

\pagebreak

\subsection{Solutions}
\paragraph{Scale a row}
\begin{minted}[bgcolor=mygreen!10]{c}
  void row_scale(int n, float R[n], float fac) {
    for (int j = 0; j < n; j = j + 1) {
      R[j] = fac * R[j];
    }
    return;
  }
\end{minted}

\paragraph{Add a row}
\begin{minted}[bgcolor=mygreen!10]{c}
  void row_add(int n, float R0[n], float R1[n], float fac) {
    for (int j = 0; j < n; j = j + 1) {
      R0[j] = R0[j] + fac * R1[j];
    }
    return;
  }
\end{minted}

\paragraph{Swap a row}
\begin{minted}[bgcolor=mygreen!10]{c}
  void row_swap(int n, float R0[n], float R1[n]) {
    float tmp = 0.0;
    for (int j = 0; j < n; j = j + 1) {
      tmp = R0[j];
      R0[j] = R1[j];
      R1[j] = tmp;
    }
    return;
  }
\end{minted}

\paragraph{Print a matrix}
\begin{minted}[bgcolor=mygreen!10]{c}
  #include <stdio.h>

  void matrix_print(int m, int n, float M[m][n]) {
    for (int i = 0; i < m; i = i + 1) {
      printf("[");
      for (int j = 0; j < n - 1; j = j + 1) {
        printf("%.2f ", M[i][j]);
      }
      printf("%.2f]\n", M[i][n-1]);
    }
    printf("\n");
    return;
  }
\end{minted}

\pagebreak

\subsection{Row reduction algorithm}
The row reduction (also know as gaussian elimination) algorithm, is used to
modify the form of a matrix into upper triangular without modifying its
kernel. This algorithm is used to solve linear systems of equations.
\begin{figure}[h]
  \centering
  \includegraphics[width=\textwidth]{flow-diagram-informal.png}
  \caption{
    This is the (informal) flow diagram of the row-reduction
    algorithm for matrices.
  }
  \label{fig:flow-diagram-informal}
\end{figure}

\paragraph{Implementation}
\begin{minted}[bgcolor=mygreen!10]{c}
  int main(void) {
    float M[3][4] = {
      {0, 2, 3, 4},
      {1, 7, 7, 7},
      {9, 8, 7, 0}
    };
    row_reduction(3, 4, M);
    matrix_print(3, 4, M);
    return 0;
  }
\end{minted}

\begin{minted}[bgcolor=mygreen!10]{c}
  void row_reduction(int m, int n, float M[m][n]) {
    int oi = 0;
    int oj = 0;
    main:
    if (M[oi][oj] != 0) {
      reduction:
      row_scale(n, M, 1/M[oi][oj]);
      for (int i = oi + 1; i < m; i = i + 1) {
        row_sub(n, M[i], M[oi], M[i][oj]);
      }
      if (oi < m) {
        oi = oi + 1;
        goto main;
      }
      else {
        return;
      }
    }
    else {
      if (exists_row_nonzero(m, n, M)) {
        row_swap_nearest_zero(m, n, M);
        goto reduction;
      }
      else {
        if (oj < n) {
          oj = oj + 1;
          goto main;
        }
        else {
          return;
        }
      }
    }
  }
\end{minted}


\pagebreak

\paragraph{Optimizations}
After going through the hastle of translating the abstract algorithm
represented by a diagram into the C language syntax and verifying
its expected behaviour, one might be left unsatisfied and have some
of these questions in mind:
\begin{itemize}
\item
  \textit{The diagram is very simple, the code is very ugly. But it works
  great without any problems so its fine right?} It is until it isn't.
  Computer science tells us that there is no way of catching certain
  errors prior to running the program. Ergo write everything with the
  intention of making it easier to fix when it will break.
\item
  \textit{What has been just done can be applied to any language. Is the only
  advantage of C that of making it faster?}
  Although the fastest programs are written in C, that doesn't imply that
  writing any program in C will make it fast.
  The sole dictator of a program's speed is the hardware and the algorithm.
  The race between languages is a race where everyone gets the same car and
  are is only allowed to remove things from it. The C language doesn't care
  about its driver's well being (you), it just wants to win the race:
  no carpets, no seats, no airbags, no safety, no spedometer...
  If you are the best driver you will definitely win the race, otherwise
  you will just die (the program won't run and that will make you want
  to be not so nice towards yourself and other).
\item
  \textit{If my focus isn't performance is it worth learning C?}
  Yes it is. Following the previous answer, it has no other structure other
  than the absolutely necessary one.
\end{itemize}

In the implementation diagram \ref{fig:flow-diagram-formal} pay attention to
the row swapping section. There is obviously some useless repetition carrying
over from checking the possibility of doing a row swap at all, more importantly
it introduces a hanging terminating condition (which shouldn't resolve but is
still considered bad practice). It is better to remove it all together and
adjust it so that the improved \ref{fig:flow-diagram-formal-1}

\begin{figure}[p]
  \centering
  \includegraphics[height=.9\textheight]{row-reduction-diagram-formal.png}
  \caption{Formalized diagram of the previous row reduction algorithm implementation.}
  \label{fig:flow-diagram-formal}
\end{figure}

\begin{figure}[p]
  \centering
  \includegraphics[height=.9\textheight]{row-reduction-diagram-formal-1.png}
  \caption{Formalized diagram of the previous row reduction algorithm implementation.}
  \label{fig:flow-diagram-formal-1}
\end{figure}

\begin{figure}[p]
  \centering
  \includegraphics[height=.9\textheight]{row-reduction-diagram-formal-2.png}
  \caption{Formalized diagram of the previous row reduction algorithm implementation.}
  \label{fig:flow-diagram-formal-2}
\end{figure}

\begin{figure}[p]
  \centering
  \includegraphics[height=.9\textheight]{row-reduction-diagram-formal-3.png}
  \caption{Formalized diagram of the previous row reduction algorithm implementation.}
  \label{fig:flow-diagram-formal-3}
\end{figure}

\begin{figure}[p]
  \centering
  \includegraphics[height=.9\textheight]{row-reduction-diagram-formal-4.png}
  \caption{Formalized diagram of the previous row reduction algorithm implementation.}
  \label{fig:flow-diagram-formal-4}
\end{figure}

\begin{figure}[p]
  \centering
  \includegraphics[height=.9\textheight]{row-reduction-diagram-formal-5.png}
  \caption{Formalized diagram of the previous row reduction algorithm implementation.}
  \label{fig:flow-diagram-formal-5}
\end{figure}

\pagebreak

\begin{minted}[bgcolor=mygreen!10]{c}
  void row_reduction(const int m, const int n, float M[m][n]) {
    for (
      int oi = 0, oj = 0, pi = oi;
      (oi < m && oj < n);
      {oi++; oj++; pi = oi;}
    ) {
      // Find pivot row.
      while (M[pi][oj] == 0.0) {
        if (pi < m - 1) {
          pi++;
        }
        else if (oj < n - 1) {
          pi = oi;
          oj++;
        }
        else {
          return;
        }
      }
      // Normalize pivot row.
      row_scale(n - oj, M[pi] + oj, 1/M[pi][oj]);
      // Clip offset column using pivot row as a mask.
      for (int i = pi + 1; i < m; i++) {
        row_add(n - oj, M[i] + oj, M[pi] + oj, -M[i][oj]);
      }
      // Swap pivot row with offset row.
      if (pi != oi) {
        row_swap(n - oj, M[oi] + oj, M[pi] + oj);
      }
    }
    return;
    /*
       [NOTES]
       --------------------------------------
       Everything preceding the offset during
       matrix triangularization is such that:
       1) It doesn't change further steps.
       2) It isn't changed by further steps.
       --------------------------------------
       therefore it is completely ignored by
       offsetting the matrices index.
       --------------------------------------
    */
  }
\end{minted}


\pagebreak

\section{Data representation and ordering}
\subsection{Scope and pointers}
\paragraph{FAQ}
\begin{itemize}
\item What are pointers? An integer type.
\item What do pointers do? I don't know, whatever you want your integers to do.
\item What uses pointers? Referencing and de-referencing operators.
\item What does (de-)referencing do? Allows you to access parallel scopes.
\item What is a scope? A device mounted on top of a rifle.
\end{itemize}
\paragraph{Scope}
Up until this point we overlooked the idea of scope.
In the C-language a scope is denoted by the curly braces enclosure.

\begin{minted}[bgcolor=mygreen!10]{c}
  {
    int a = 10;
    int a = 23;
    printf("%d\n", a);
    // error there are to many 'a'
  }

  {
    int a = 10;
    {
      int a = 23;
      printf("%d\n", a);
      // shows 23 on terminal
    }
    printf("%d\n", a);
    // shows 10 on terminal
  }

  {
    int a = 10;
    {
      int b = 67;
      printf("%d\n", a);
      // shows 10 on terminal
    }
    printf("%d\n", b);
    // error there is no such b
  }
\end{minted}

What happens when I call a function? Why does this happen?
\begin{minted}[bgcolor=mygreen!10]{c}
  void funny(int a) {
    a = 67;
    return;
  }

  int main(void) {
    int a = 10;
    funny(a);
    printf("%d\n", a);
    // shows 10 on terminal
    return 0;
  }
\end{minted}
Lets observe the scope, this is what happens when we call a function
\begin{minted}[bgcolor=mygreen!10]{c}
  funny:
  {
    int a = 10;
    a = 67;
    goto _return;
  }

  {
    int a = 10;
    goto funny;
    return:
    printf("%d\n", a);
  }
\end{minted}

What about a pointer?
What happens when I call a function? Why does this happen?
\begin{minted}[bgcolor=mygreen!10]{c}
  void funny(int* a) {
    *a = 67;
    return;
  }

  int main(void) {
    int a = 10;
    funny(&a);
    printf("%d\n", a);
    // shows 10 on terminal
    return 0;
  }
\end{minted}
\begin{minted}[bgcolor=mygreen!10]{c}
  funny:
  {
    int* a = &a;
    *a = 67;
    goto _return;
  }

  {
    int a = 10;
    goto funny;
    return:
    printf("%d\n", 67);
  }
\end{minted}

Couldn't I just return 'a' instead? Yes but what if you want to change more than one variable?


\subsection{Arrays}
\subsection{Integers}
\subsection{Decimals}
\subsection{Fixed point and floating point}


\pagebreak

\section{Biblio}
\begin{itemize}
\item ``Informatica: Fondamenti Culturali e Tecnologici 1986'' by \textit{Michele Pellerey}
\item ``Introduction to Algorithms (third edition) 2009'' by \textit{Thomas H. Cormen}
\end{itemize}

\end{document}
